% ===== §2 =====
\section{A Taxonomy of Definitions of Culture}
\label{p24:sec:taxonomy}

\noindent
Before the several theories of culture are taken up one by one, they should be laid out as a field, for they do not merely differ in emphasis but disagree at the root, in what they take culture \emph{to be}. A theory of cultural formation always rests on a prior commitment about the mode of existence of its object, and it is at that level, and not at the level of particular findings, that the position of this paper is to be located. The purpose of this section is therefore cartographic. It sets two axes across which the major commitments can be plotted, places each of the traditions the following sections examine at its coordinate, and shows that the position this paper will occupy is not one more entry in a list but a region of the map the existing traditions have left vacant. To argue for the originality of a view by exhibiting an unoccupied coordinate is more honest than to assert it, because it makes the claim answerable: if the coordinate were already filled, the argument would fail, and it can be seen at a glance whether it is.

\subsection*{Two axes}

The first axis concerns the \emph{mode of being} of culture, and it runs from the settled to the generative. At one pole culture is a completed and transmissible entity, a heritage that exists in the manner of an object, handed from one bearer to the next and in principle the same across the handing. At the other pole culture is an ongoing process, unfinished by nature, existing only in the manner of an event that is always still occurring. This is the axis on which the etymological tension of the previous section, between the tended possession and the ongoing transformation, becomes a theoretical variable.

The second axis concerns the \emph{bearer} of culture, and it runs from the individual to the relational-or-structural. At one pole the individual is the unit that carries culture, its container or its vehicle, and culture is what is deposited in persons. At the other pole culture obtains between or above persons, in a structure or a relation that no single individual instances, so that the individual is an effect or a moment of the culture rather than its container. These two axes are not independent in practice, but they are distinct in principle, and plotting the traditions against both at once discloses a structure that neither axis alone would show.

\subsection*{The map}

Plotted against these axes, the traditions examined in the following sections fall as the table records. Each row states, besides the coordinate, what the tradition takes culture to be, how it treats drive, and how it treats power, so that the two lines this paper carries through its ontology---the psychoanalytic and the question of power---can be seen entering, or failing to enter, each account.

\begin{center}
\small
\renewcommand{\arraystretch}{1.45}
\begin{tabularx}{\linewidth}{@{}>{\RaggedRight\arraybackslash}p{2.5cm} >{\RaggedRight\arraybackslash}p{2.7cm} >{\RaggedRight\arraybackslash}X >{\RaggedRight\arraybackslash}p{2.5cm} >{\RaggedRight\arraybackslash}p{2.5cm}@{}}
\toprule
\textcolor{rosedeep}{\textbf{Tradition}} & \textcolor{rosedeep}{\textbf{Coordinate}} & \textcolor{rosedeep}{\textbf{What culture is}} & \textcolor{rosedeep}{\textbf{Drive}} & \textcolor{rosedeep}{\textbf{Power}} \\
\midrule
Transmission \& collective representation \newline\emph{\small Tylor, Durkheim} & settled \texttimes{} collective & a learned, transmitted whole; collective representation, social fact & untouched; culture as outer norm & unthematized; power hidden in the coercion of the ``social fact'' \\
Symbol \& interpretation \newline\emph{\small Geertz, Cassirer} & settled \texttimes{} structural & a prior web of meaning, a system of symbolic forms & symbolic order as big Other; the subject meets lack on entering language & the power dimension of meaning left undeveloped \\
Structure \& reproduction \newline\emph{\small L\'evi-Strauss, Bourdieu} & settled-tending-dynamic \texttimes{} structural & deep structure; cultural capital, accruable, convertible, extractable & desire organized by the logic of distinction, its drive-root undug & \emph{central}: culture as the hidden mechanism of class reproduction \\
Power as ontology \newline\emph{\small Gramsci, Foucault, Hall} & dynamic \texttimes{} structural (power congealed) & the encoding of a power order; no culture outside power & inside the symbolic via the phallic signifier (the weld) & \emph{ontological}: power structure defines culture structure \\
The dialogic \newline\emph{\small Bakhtin} & generative \texttimes{} dialogic & the unfinished, generated in answering & answerability touches the other's desire, unsystematized & heteroglossia implies a counter-hegemony, untheorized \\
The Chinese \emph{wen--hua} \newline\emph{\small Yijing, Confucian, Daoist} & generative \texttimes{} civilizational & \emph{hua}: transformation, bringing-to-form; culture as ongoing act & sublimation implicit in \emph{liyue} (rites and music), undertheorized & pedagogy from above; the power face of \emph{huacheng} \\
\midrule
\textcolor{rosedeep}{\textbf{GRB (this paper)}} & \textbf{generative \texttimes{} relational (vacant)} & the dialectical exteriorization, across three registers, of interwoven drive and generative history & \emph{ground}: culture as coding, sublimation, the sediment of shared fantasy & power is internal to culture, yet just generation generates the \emph{counter-power} against its entrenchment \\
\bottomrule
\end{tabularx}
\end{center}

\subsection*{Reading the map}

Read downward, the table shows a drift rather than a scatter. The mode of being moves from the settled toward the generative; the bearer moves from the individual and the structural toward the relational; and the two columns on the right, drive and power, move from the untouched toward the thematic, so that what the earliest traditions do not register at all becomes, in the lower rows, the explicit matter of the account. The lowest row, the position of this paper, occupies the coordinate the drift approaches but none of the existing traditions reaches: the generative and the relational together, with drive taken as the ground of culture and power taken as internal to it yet answerable, within the same generative movement, to a counter-power that prevents its entrenchment. That the phallic signifier appears in the drive column of the power tradition and the counter-power in the power column of this paper's own row is not incidental. It marks the two welds by which the psychoanalytic line and the line of power are joined to each other and to the whole, and both are redeemed in their places below.

Two uses follow, and they set the plan of the rest of Part One. The task of the antithesis, when it comes, will be to draw the upper six traditions from their several coordinates toward the vacant one, not by refuting each in turn but by exposing the presupposition they share and inverting it. The task of the synthesis will be to name and define the coordinate so reached. The sections that immediately follow, however, owe each tradition a fair hearing on its own terms before any of this can be earned, and they proceed in the order of the table, from the tradition that holds culture most nearly as a settled thing to the two that already lean, half a foot, toward the generation this paper will claim.
